\documentclass[aspectratio=169,10pt,spanish]{beamer}

\input{../../../_preambulo_beamer.tex}
\input{../../../_comandos_beamer.tex}

\selectlanguage{spanish}

\title{MA1001B - Modelación Estadística para la Toma de Decisiones}
\subtitle{Lección 36: Prueba de hipótesis para una proporción de negocio}
\author{}
\institute{Tecnol\'ogico de Monterrey}
\date{}

\begin{document}

\begin{frame}[plain]
  \vspace{1.2cm}
  {\Large\bfseries MA1001B - Modelación Estadística para la Toma de Decisiones\par}
  \vspace{0.7cm}
  {\large Lección 36: Prueba de hipótesis para una proporción de negocio\par}
  \vspace{0.7cm}
  {\color{TecRojo}\rule{\textwidth}{1.2pt}}
  \vspace{0.8cm}

  Facultad MA1001B\\[0.3cm]
  Periodo\\[0.3cm]
  Tecnol\'ogico de Monterrey
\end{frame}

\begin{frame}{La pregunta de una proporción}
  \begin{alertblock}{Pregunta guía}
    ¿Cómo se prueba una tasa reclamada de banderas de retraso $p=0.10$ a
    partir de $x=28$ banderas en $n=200$ lotes?
  \end{alertblock}

  \vspace{0.6em}
  Al final de esta lección, usted puede:
  \begin{itemize}
    \item escribir $H_0:p=0.10$ frente a $H_1:p>0.10$;
    \item comprobar $np_0>5$ y $n(1-p_0)>5$;
    \item calcular $z$ con $p_0$ en el error estándar;
    \item explicar por qué reemplazar $p_0$ por $\hat{p}$ es el estadístico
      incorrecto.
  \end{itemize}

  \vfill
  \scriptsize Todos los lotes y banderas de retraso usados hoy son sintéticos. No se usan datos reales de empresa.
\end{frame}

\begin{frame}{Una tasa reclamada y un conteo muestral}
  \small
  Un proceso sintético de inspección de entrada afirma que el 10 por ciento
  de los lotes se marcan con retraso. El registro observado es $x=28$
  banderas de retraso en $n=200$ lotes.

  \[
    \hat{p}=\frac{28}{200}=0.14,\qquad
    np_0=200\times0.10=20,\qquad
    n(1-p_0)=180.
  \]

  \begin{exampleblock}{Hipótesis de cola derecha}
    $H_0:p=0.10$ frente a $H_1:p>0.10$, con $\alpha=0.05$. La aproximación
    normal $z$ se usa solo después de $np_0>5$ y $n(1-p_0)>5$.
  \end{exampleblock}
\end{frame}

\begin{frame}[fragile]{Paso 1: Proporción muestral y condiciones}
  \begin{columns}[T]
    \begin{column}{0.42\textwidth}
      \small
      \begin{lstlisting}
phat = 28 / 200
np0 = n * p0
nq0 = n * (1 - p0)
      \end{lstlisting}

      \begin{block}{Salida verificada}
        $\hat{p}=0.140000$\\
        $np_0=20$, $n(1-p_0)=180$\\
        Condiciones cumplidas: sí
      \end{block}
    \end{column}
    \begin{column}{0.58\textwidth}
      \centering
      \includegraphics[width=\textwidth]{../figuras/prueba_proporcion_01_muestra_y_condiciones.png}
      \vspace{0.2em}
      \scriptsize $p_0$ hipotética frente a $\hat{p}$ del Paso 1
    \end{column}
  \end{columns}
\end{frame}

\begin{frame}{El error estándar de la prueba usa $p_0$}
  \small
  Bajo $H_0:p=p_0$ la varianza muestral es $p_0(1-p_0)/n$, no
  $\hat{p}(1-\hat{p})/n$:
  \[
    z=\frac{\hat{p}-p_0}{\sqrt{p_0(1-p_0)/n}}
    =\frac{0.14-0.10}{0.021213}=1.885618.
  \]
  El valor p de cola derecha es
  \[
    p=P(Z\ge 1.885618)=0.029673.
  \]
  Como $p<\alpha=0.05$ y $z>z^*=1.644854$, la prueba rechaza $H_0$.
\end{frame}

\begin{frame}[fragile]{Paso 2: $z$ y el valor p de cola derecha}
  \begin{columns}[T]
    \begin{column}{0.40\textwidth}
      \small
      \begin{lstlisting}
se = sqrt(p0*(1-p0)/n)
z = (phat - p0) / se
p = norm.sf(z)
      \end{lstlisting}

      \begin{block}{Salida verificada}
        $EE=0.021213$, $z=1.885618$\\
        $p=0.029673$, $z^*=1.644854$\\
        Decisión: rechazar $H_0$
      \end{block}
    \end{column}
    \begin{column}{0.60\textwidth}
      \centering
      \includegraphics[width=\textwidth]{../figuras/prueba_proporcion_02_z_valor_p.png}
      \vspace{0.2em}
      \scriptsize Prueba $z$ de cola derecha del Paso 2
    \end{column}
  \end{columns}
\end{frame}

\begin{frame}{Paso 3: $\hat{p}$ en el EE es el estadístico incorrecto}
  \begin{columns}[T]
    \begin{column}{0.42\textwidth}
      \small
      El EE correcto usa $p_0$: $0.021213$, $z=1.885618$, $p=0.029673$,
      rechazar.

      \vspace{0.5em}
      El EE incorrecto usa $\hat{p}$: $0.024536$, $z=1.630278$, $p=0.051521$,
      no rechazar.

      \vspace{0.5em}
      \textbf{Límite:} una prueba de $H_0:p=p_0$ debe poner $p_0$ en el error
      estándar. Usar $\hat{p}$ es un EE de intervalo de confianza y puede
      invertir la decisión con $\alpha=0.05$.
    \end{column}
    \begin{column}{0.58\textwidth}
      \centering
      \includegraphics[width=\textwidth]{../figuras/prueba_proporcion_03_limite_ee_incorrecto.png}
      \vspace{0.2em}
      \scriptsize EE correcto frente a EE incorrecto del Paso 3
    \end{column}
  \end{columns}
\end{frame}

\begin{frame}{Verificación de conceptos}
  \begin{enumerate}
    \item Calcule $\hat{p}$ a partir de $x=28$ y $n=200$.
    \item ¿Por qué se comprueba $np_0=20$ con $p_0=0.10$ y no con $\hat{p}$?
    \item Escriba la fórmula $z$ que pertenece a esta prueba.
    \item ¿Por qué usar $\hat{p}$ en el EE puede invertir la decisión con
      $\alpha=0.05$?
  \end{enumerate}

  \vspace{0.6em}
  \begin{block}{Conexión con la Lección 37}
    La sesión puede continuar directamente con una prueba $t$ de dos muestras
    para dos medias independientes, donde cada grupo tiene su propia $s$ y
    $n$.
  \end{block}

  \vfill
  \scriptsize Fuente: OpenStax, \textit{Introductory Business Statistics 2e},
  Capítulo 9, Secciones 9.3--9.4. CC BY-NC-SA.
\end{frame}

\appendix



\end{document}
